% TOP_PROBLEM: {"id": 20001046, "title": "An arc--nonedge defect identity for Sullivan-1", "category_id": 2, "category": {"id": 2, "name": "combinatorics", "display_name": "Combinatorics", "description": "Counting problems, graph theory, discrete structures.", "slug": "combinatorics", "order_index": 2, "created_at": "2026-07-31T15:26:25.671Z"}, "status": "partially_solved", "problem_number": "AIM-COMBINATORICS-0171", "set_id": 14}
% FIRST_POSTED: 2026-10-01
% ENTRY_KIND: statement_only
% UnsolvedMath ID 20001046; source catalogue code AIM-COMBINATORICS-0171
% !TeX program = lualatex
% Definition and sources only; this file is not an LLM solution attempt.
\documentclass[11pt,a4paper]{article}
\usepackage[margin=25mm]{geometry}
\usepackage{fontspec}
\setmainfont{lmroman10-regular.otf}[BoldFont=lmroman10-bold.otf,
 ItalicFont=lmroman10-italic.otf,BoldItalicFont=lmroman10-bolditalic.otf]
\usepackage{amsmath,amssymb,mathtools}
\usepackage{unicode-math}
\setmathfont{latinmodern-math.otf}
\AtBeginDocument{\let\setminus\smallsetminus}
\usepackage{xurl,hyperref}
\hypersetup{colorlinks=true,urlcolor=blue}
\setcounter{secnumdepth}{0}
\setlength{\parindent}{0pt}
\setlength{\parskip}{5pt}
\setlength{\emergencystretch}{3em}
\begin{document}
\section*{An arc--nonedge defect identity for Sullivan-1}\label{20001046}
{\small UnsolvedMath ID 20001046 · AIM-COMBINATORICS-0171 · Combinatorics · AIM Workshop Problem Lists}
\subsection{Definitions and mathematical statement}
Consider a nonempty finite oriented graph $D=(V,E)$: there are no
loops, no parallel arcs, and no pair of oppositely directed arcs.
For $v\in V$, its in-neighbourhood is $N^-(v)=\{u:u\to v\}$ and
its first out-neighbourhood is $N^+(v)=\{u:v\to u\}$.
Its second out-neighbourhood is
\[
 N_2^+(v)=\left(\bigcup_{u\in N^+(v)}N^+(u)\right)
                         \setminus\bigl(N^+(v)\cup\{v\}\bigr).
\]
Thus its members are at shortest directed-path distance exactly two.
The union eliminates multiple routes to an endpoint, and the
subtraction eliminates immediate out-neighbours.

The compromise conjecture, Conjecture 6.6 in the AIM source, asks
whether every such graph has a vertex $v$ with
\[
                              |N_2^+(v)|\ge |N^-(v)|.
\]
The comparison is with incoming neighbours, while the two-step
reachability is in the outgoing direction. Reversing this direction
in just one of the two definitions would change the conjecture.
The witnessing vertex can depend on $D$, and neither connectivity
nor degree regularity is assumed. Sources, having indegree zero,
are immediate witnesses, so a counterexample would have positive
indegree at every vertex. Opposite arcs remain forbidden even in
that case. The target counts distinct second neighbours, not the
number of directed two-edge paths or the number of vertices at
out-distance at most two.
\subsection{Short English statement}
Must every oriented graph have a vertex with at least as many second out-neighbours as incoming neighbours?
\subsection{Sources}
\begin{itemize}
\item[S1] ULAM.AI and the UnsolvedMath Contributors, supplied problems.json, record 20001046 (AIM-COMBINATORICS-0171), AIM Workshop Problem Lists. Catalogue identity and source transcription; CC BY 4.0.
\url{https://huggingface.co/datasets/ulamai/UnsolvedMath}.
\item[S2] American Institute of Mathematics, The Caccetta-Haggkvist conjecture, item 6.6. Original workshop question underlying AIM-COMBINATORICS-0171.
\url{https://aimath.org/WWN/caccetta/caccetta.pdf}.
\end{itemize}
\end{document}
